% ============================================================================
% Chapter 4 --- Institutional self-assessment
% ----------------------------------------------------------------------------
% Purpose : provide an auto-questionnaire allowing any institution to
%           position itself before reading Parts II and III. Replaces the
%           original A.4 "open questions" with an instrument turned toward
%           the institution rather than toward the design.
% Page target: ~8 pages.
% Dependencies: applies the dimensions defined in Ch. 3 (sovereignty grid)
%               and the scoring scale defined there.
% ----------------------------------------------------------------------------
% Decisions registered in _coherence-EN.md:
%   D4.1 self-assessment is conducted by the institution itself (no external audit)
%   D4.2 four diagnostic sections (A--D) plus one positioning section (E)
%   D4.3 questions phrased to be verifiable, compact, diagnostic, non-leading
%   D4.4 trajectory selection matrix gives heuristics, not algorithmic decision
% ============================================================================

\chapter{Institutional self-assessment}
\label{ch:self-assessment}

\begin{caveatbox}[Dependencies]
\noindent This chapter applies the five-dimension grid and the
four-point scoring scale established in
\trchap{ch:sovereignty-grid}{The sovereignty grid}. Section~E
(\S\ref{sec:sa-sovereignty}) presupposes that the reader has
completed Chapter~3 of the Technical Reference. The trajectory
selection matrix (\S\ref{sec:sa-results}) refers forward to
\trchap{ch:playbook-reading}{How to read this playbook} for the
definitions of the three migration trajectories.
\end{caveatbox}

The instrument presented in this chapter is a self-administered
diagnostic. Its purpose is not to grade an institution against an
external benchmark but to make explicit the starting position from
which any migration toward the architecture and playbook of this
Reference would proceed. An honest self-assessment, conducted by the
institution itself with the aid of its own documentation, is more
useful than an externally produced report whose assumptions the
institution has not internalised.

%-----------------------------------------------------------------------------
\section{Purpose and method}
\label{sec:sa-purpose}

The self-assessment is composed of five sections. Sections~A--D are
diagnostic: they characterise the institution's current technical
maturity along the four planes that the migration playbook will
exercise (identity, network, observability, endpoint). Section~E is
positioning: it applies the sovereignty grid of
\trchap{ch:sovereignty-grid}{} to the institution as a whole, asking
both \emph{where it stands today} and \emph{where it aims to stand}.

\paragraph{Time and team.} An indicative envelope for completing the
assessment is a single facilitated session of two to three hours
within a small cross-functional group --- typically one identity
engineer, one network engineer, one security or governance
representative, and one person familiar with procurement. Institutions
will adjust the envelope to their own deliberation cadence; the
resulting profile is a working artefact for the institution, not a
deliverable for an external party.

\paragraph{Scoring.} Each question admits one of four answers, using
the same four-point scale as the sovereignty grid: \textbf{0 Absent},
\textbf{1 Partial}, \textbf{2 Established}, \textbf{3 Robust}. Half-
points are not allowed; when uncertainty exists, the lower score is
recorded with a note on the basis for uncertainty. The most useful
output is not the aggregate but the distribution: a section in which
most answers are 1 or 2 with one outlier 0 reveals more than an
average score would.

\paragraph{Honesty discipline.} The instrument is only useful if the
institution scores itself as it actually is, not as it would prefer to
appear. Two practices help: scoring is performed before the trajectory
matrix (\S\ref{sec:sa-results}) is consulted, so respondents do not
calibrate to a preferred outcome; and the basis for each non-zero
score is briefly documented (one line) so that the team can revisit
its judgement later.

%-----------------------------------------------------------------------------
\section{Section A --- Identity and federation maturity}
\label{sec:sa-identity}

This section diagnoses the maturity of the institution's identity
plane --- the layer that the first migration wave
(\trchap{ch:wave1-identity}{Wave~1 --- Identity and federation}) will
exercise.

\begin{enumerate}[leftmargin=2em,nosep,label=A\arabic*.]
  \item Is there a single canonical identity provider for
  \loc{institution-name}, or multiple competing ones?

  \item Does the institution participate in \gls{edugain} (or a
  regional equivalent) for inter-institutional federation?

  \item Is account lifecycle (creation, role change, departure)
  automated end-to-end via \gls{scim} or equivalent?

  \item Is \gls{mfa} enforced for all staff and student accounts, with
  documented exceptions?

  \item Are role and group memberships expressed in a single source of
  truth, or duplicated across systems?

  \item Is the identity provider's source code or licence under
  institutional control (open source, or contractual right to operate
  independently)?

  \item Are session lifetimes and elevation policies documented and
  consistent across applications?

  \item Is there a documented break-glass procedure for emergency
  access that does not bypass audit?

  \item Is the identity provider's evolution roadmap controlled by the
  institution, by an external supplier, or by community consensus?

  \item Has the identity provider been independently audited within
  the last \loc{assurance-recency-window}?
\end{enumerate}

%-----------------------------------------------------------------------------
\section{Section B --- Network and segmentation maturity}
\label{sec:sa-network}

This section diagnoses the maturity of the network plane that
\trchap{ch:wave3-network}{Wave~3 --- Network segmentation} will
exercise.

\begin{enumerate}[leftmargin=2em,nosep,label=B\arabic*.]
  \item Are network zones (administrative, research, teaching,
  \gls{byod}, \gls{iot} where applicable, sensitive) declared and
  documented?

  \item Does each zone have a published security profile (filtering
  rules, egress policy, audit obligations)?

  \item Is there a \gls{sciencedmz} for high-throughput research
  traffic, or an explicit architectural decision against one?

  \item Is east-west traffic between zones restricted by default, or
  permitted by default?

  \item Are \gls{nac} policies enforced at the wired and wireless
  edge?

  \item Is encryption-in-transit enforced for institutional traffic
  (\gls{tls}-everywhere posture)?

  \item Are flow records (NetFlow, IPFIX, sFlow) collected and
  retained for a documented period?

  \item Is there a documented \gls{byod} segregation strategy
  (separate \gls{vlan}, captive portal, posture verification)?

  \item Are network device configurations under version control with
  reviewable history?

  \item Has the segmentation been independently penetration-tested
  within the last \loc{assurance-recency-window}?
\end{enumerate}

%-----------------------------------------------------------------------------
\section{Section C --- Observability and audit maturity}
\label{sec:sa-observability}

This section diagnoses the maturity of the observability plane that
\trchap{ch:wave2-observability}{Wave~2 --- Observability baseline}
will exercise and that \trchap{ch:wave7-audit-ir}{Wave~7 --- TAC-aware
audit and IR} will refine.

\begin{enumerate}[leftmargin=2em,nosep,label=C\arabic*.]
  \item Is there a central log collection point (\gls{siem} or
  equivalent) covering at least identity, network, and endpoint?

  \item Are log sources catalogued and their criticality documented?

  \item Is log retention specified per data classification and
  contractually enforceable?

  \item Are detection rules (correlations, alerts) documented and
  version-controlled?

  \item Is there a continuous monitoring capacity, internal or
  contracted, with defined response times?

  \item Are incident response runbooks documented and exercised?

  \item Is the \gls{siem} cost predictable (fixed envelope, not
  per-GB ingestion exposed to volume spikes)?

  \item Are log schemas standardised (OpenTelemetry, CloudEvents, or
  an institutional canonical schema)?

  \item Has an incident response drill been conducted within the last
  \loc{assurance-recency-window}?

  \item Is there a documented procedure to preserve forensic evidence
  during an incident?
\end{enumerate}

%-----------------------------------------------------------------------------
\section{Section D --- Endpoint management maturity}
\label{sec:sa-endpoint}

This section diagnoses the maturity of the endpoint plane that
\trchap{ch:wave6-endpoint}{Wave~6 --- Endpoint and posture} will
exercise.

\begin{enumerate}[leftmargin=2em,nosep,label=D\arabic*.]
  \item Is the endpoint estate inventoried and classified (managed,
  \gls{byod}, \gls{iot})?

  \item Is mobile-device management coverage in place for each
  managed-endpoint operating system in use?

  \item Are posture verification policies documented and reviewed?

  \item Is patch service-level documented and measured against
  reality?

  \item Is \gls{byod} acceptable use enforced technically, or relied
  upon contractually only?

  \item Are administrative privileges on workstations centrally
  controlled with a documented exception process?

  \item Has the endpoint-management tooling been evaluated against the
  sovereignty grid (sources, exit assistance, data residency)?

  \item Is there a documented decommissioning procedure (data wipe,
  account deactivation, return of asset)?
\end{enumerate}

%-----------------------------------------------------------------------------
\section{Section E --- Sovereignty positioning}
\label{sec:sa-sovereignty}

This section applies the sovereignty grid of
\trchap{ch:sovereignty-grid}{} to the institution as a whole. For each
of the five dimensions, two questions are asked: one diagnostic of
the current state, one diagnostic of the institution's stated intent.
A gap between the two is itself useful information about the maturity
of the governance.

\paragraph{E1 --- Data sovereignty.}
\begin{enumerate}[leftmargin=2em,nosep,label=E1.\arabic*.]
  \item Are storage location, key custody, and disclosure regime
  documented per data class?
  \item Is \emph{``data dimension Established or better''} set as a
  procurement gate for sensitive workloads, in writing?
\end{enumerate}

\paragraph{E2 --- Technical sovereignty.}
\begin{enumerate}[leftmargin=2em,nosep,label=E2.\arabic*.]
  \item Can the institution inspect or independently rebuild every
  critical component of the infrastructure?
  \item Is replaceability written into \gls{rfp} language as a
  structural requirement, not as a desirable preference?
\end{enumerate}

\paragraph{E3 --- Financial sovereignty.}
\begin{enumerate}[leftmargin=2em,nosep,label=E3.\arabic*.]
  \item Is multi-year operating cost for the digital infrastructure
  predictable and capped contractually?
  \item Is supplier substitutability evaluated annually as a
  sovereignty exercise, independent of incident response?
\end{enumerate}

\paragraph{E4 --- Operational sovereignty.}
\begin{enumerate}[leftmargin=2em,nosep,label=E4.\arabic*.]
  \item Are the skills for each critical layer present in the
  institution or accessible through its national research and
  education network and peer network?
  \item Has mutualisation with peer institutions been considered at
  planning level for at least one shared operational burden?
\end{enumerate}

\paragraph{E5 --- Institutional sovereignty.}
\begin{enumerate}[leftmargin=2em,nosep,label=E5.\arabic*.]
  \item Can the institution decline a supplier roadmap change without
  contractual penalty for the current procurement period?
  \item Is the institution publishing, or willing to publish, its
  sovereignty profile by component to its academic community?
\end{enumerate}

%-----------------------------------------------------------------------------
\section{Reading the results: trajectory selection matrix}
\label{sec:sa-results}

The matrix in Table~\ref{tab:sa-trajectory} translates a profile of
answers into a recommended dominant migration \emph{trajectory} ---
greenfield, brownfield single-suite-centric, or brownfield hybrid --- and
into an indicative \emph{pace} (5-, 7-, or 10-year horizon).

The matrix offers heuristics, not algorithmic decisions: in particular,
the trajectory and the pace are typically constrained by factors that
the questionnaire does not capture (executive mandate, budget cycle,
ongoing transformation programmes). The matrix is the entry into
\trchap{ch:playbook-reading}{}, which develops the trajectory and pace
choices in the context of the playbook itself.

\begin{table}[H]
\centering\small
\caption{Trajectory and pace heuristics from the self-assessment
profile. ``A--D maturity'' is the dominant trend across the
diagnostic sections; ``E gap'' is the gap between current and intent
in Section~E; ``Section~E intent'' is the intent score itself.}
\label{tab:sa-trajectory}
\begin{tabularx}{\textwidth}{%
  >{\hsize=0.9\hsize\linewidth=\hsize\bfseries\raggedright\arraybackslash}X%
  >{\hsize=0.9\hsize\linewidth=\hsize\raggedright\arraybackslash}X%
  >{\hsize=1.2\hsize\linewidth=\hsize\raggedright\arraybackslash}X%
}
\toprule
A--D maturity profile & E intent profile & Recommended trajectory and pace \\
\midrule
Mostly 0--1 (low across sections), no dominant supplier yet &
Intent~$\geq$~Established on at least three dimensions &
\textbf{Greenfield}, 5--7~year horizon. Rare; applies to new
institutions or to a deliberate from-scratch refoundation. \\
\addlinespace
Mostly 1--2, with a single dominant commercial productivity-and-identity
suite across identity, endpoint, and observability (e.g.\ Microsoft) &
Intent~$\geq$~Established on Data and Institutional dimensions &
\textbf{Brownfield single-suite-centric}, 7--10~year horizon. A common
starting point for institutions whose estate is anchored on a single
dominant productivity-suite incumbent; the playbook adds an open policy
plane and an open observability collector layer in front of the
existing incumbent-suite estate. \\
\addlinespace
Mostly 1--2, with mixed Linux / macOS / Windows estates and at least
one open-source identity component already in production &
Intent~$\geq$~Established on Technical and Institutional dimensions &
\textbf{Brownfield hybrid}, 5--7~year horizon. Frequent where
heterogeneous estates and at least one open-source identity component
already coexist; the playbook progressively
substitutes proprietary components at renewal moments while
preserving the open-source perimeter. \\
\addlinespace
Mostly 2--3 (high maturity across sections) &
Intent stable at current level, no significant gap &
Selective adoption: the institution may already realise much of the
target architecture. The playbook is used as a checklist, not as a
roadmap. Pace is set by the renewal cycles of remaining captive
components. \\
\addlinespace
Heterogeneous (some 0, some 3, no dominant trend) &
Any &
\textbf{Defer trajectory selection.} The heterogeneity itself is the
finding. Wave~0 (\trchap{ch:wave0-foundations}{}) is the appropriate
first step: governance and inventory before architectural commitment. \\
\bottomrule
\end{tabularx}
\end{table}

\paragraph{Pace adjustments.} The horizons indicated above are
adjusted by Section~E intent: a high intent score across all
dimensions may justify the shorter horizon; a low intent score
indicates that the longer horizon is realistic and that the
institution should not commit to faster pace at planning time.

\paragraph{Trajectory revision.} The trajectory selected at the
self-assessment is revised at the close of each wave
(\trchap{ch:reversibility}{Reversibility playbook}). It is normal for
an institution to start brownfield single-suite-centric and migrate to
brownfield hybrid as the open policy and observability planes mature;
it is uncommon, but possible, for a stalled brownfield trajectory to
be re-baselined as a greenfield refoundation in a contained perimeter.

%-----------------------------------------------------------------------------
\section{Re-assessment cadence}
\label{sec:sa-cadence}

The self-assessment is not a one-time exercise. Three cadences govern
its repetition.

\begin{description}[leftmargin=1.5em,style=nextline,nosep]
  \item[Annual review.] The full questionnaire is re-administered
    once per year, ideally in the same period as the institution's
    operational risk review. Year-on-year drift on a single question
    is more diagnostic than the absolute score in any one year.

  \item[Wave closure.] At the close of each migration wave
    (\trchap{ch:reversibility}{}), the relevant subset of questions
    is re-administered. Wave~1 (Identity) re-examines Section~A;
    Wave~3 (Network) re-examines Section~B; and so on. The grid
    questions of Section~E are re-administered after every wave.

  \item[Event-driven re-exercise.] An unscheduled re-exercise is
    triggered by any of: a significant pricing change from a
    component supplier, a supplier acquisition or restructuring, a
    court decision or regulatory change affecting the disclosure
    regime, a major outage, or a publicly disclosed supplier security
    incident.
\end{description}

The cumulative record of self-assessments forms part of the
institution's sovereignty governance archive, and is the principal
input to the international peer comparison anticipated in phase~2 of
the network of universities to which this Reference is addressed.
